% ============================================================
%  Paper 6 — Aroha Labs
%  "ACP: A Saga-Based Agent Communication Protocol with
%   Capability Schema Negotiation"
%
%  Target venues: AAMAS 2027 / ICSOC 2027 / USENIX ATC 2027
%  arXiv categories: cs.MA, cs.DC, cs.AI
% ============================================================

\documentclass[conference]{IEEEtran}
\IEEEoverridecommandlockouts

\usepackage{amsmath,amssymb,amsthm,mathtools}
\usepackage{algorithm}
\usepackage{algpseudocode}
\usepackage{graphicx}
\usepackage{booktabs}
\usepackage{multirow}
\usepackage{hyperref}
\usepackage{xcolor}
\usepackage{cite}
\usepackage{bm}
\usepackage{listings}

\newtheorem{theorem}{Theorem}
\newtheorem{lemma}[theorem]{Lemma}
\newtheorem{proposition}[theorem]{Proposition}
\newtheorem{corollary}[theorem]{Corollary}
\newtheorem{definition}{Definition}
\newtheorem{remark}{Remark}

\DeclareMathOperator{\E}{\mathbb{E}}
\newcommand{\Adv}{\mathcal{A}}
\newcommand{\Sign}{\mathsf{Sign}}
\newcommand{\Verify}{\mathsf{Verify}}
\newcommand{\DID}{\mathsf{DID}}
\newcommand{\CSN}{\mathsf{CSN}}
\newcommand{\Saga}{\mathcal{S}}
\newcommand{\Step}{\mathcal{T}}

\lstset{
  basicstyle={\ttfamily\scriptsize},
  breaklines=true,
  frame=none,
  columns=fullflexible
}

\hypersetup{
  colorlinks=true,
  linkcolor=blue!70!black,
  citecolor=green!50!black,
  urlcolor=blue!70!black
}

\begin{document}

\title{ACP: A Saga-Based Agent Communication Protocol\\
       with Capability Schema Negotiation}

\author{
  \IEEEauthorblockN{Aroha Labs Research}
  \IEEEauthorblockA{Aroha Protocol\\
    aroha-labs.com}
}

\maketitle

\begin{abstract}
We present the \emph{Aroha Communication Protocol} (ACP), the semantic
layer of the Aroha agent network stack governing how autonomous AI
agents initiate, negotiate, and complete multi-step interactions.
ACP defines seventeen message types organized into four groups—direct
invocation, negotiation, saga execution, and meta-communication—and
specifies four interaction patterns of increasing complexity: simple
request/response, streaming, negotiated, and distributed saga.  The
saga pattern adapts the classic database saga~\cite{garcia1987sagas}
to multi-agent settings, providing atomicity via last-in-first-out
(LIFO) compensating cancellations without distributed locking.  ACP
couples the saga lifecycle with a three-path \emph{Capability Schema
Negotiation} (CSN) mechanism that resolves capability mismatches via
cache lookup, structural JSON Schema matching~\cite{w3cjsonschema}, or
LLM-assisted intent matching, with subsequent interactions benefiting
from sub-millisecond cache hits.  A pluggable human-in-the-loop (HITL)
approval gate sits between the reserve and commit phases of each saga
step, enabling human oversight at granular transactional boundaries.
Pluggable settlement adapters integrate with Stripe, SQL ledgers, and
on-chain escrow.  We prove three formal properties: saga completeness,
ArohaCancel idempotency, and CSN cache correctness.  ACP complements
the companion protocol-stack paper~\cite{aroha2025p5} and the
reputation system of Paper~1~\cite{aroha2025p1}.
\end{abstract}

\begin{IEEEkeywords}
multi-agent systems, saga pattern, distributed transactions, capability
negotiation, human-in-the-loop, agent communication protocol
\end{IEEEkeywords}

% ============================================================
\section{Introduction}
\label{sec:intro}
% ============================================================

Autonomous AI agents increasingly act on behalf of humans in
open, heterogeneous networks~\cite{guo2024agents}: booking travel,
executing financial transactions, orchestrating chains of specialized
sub-agents.  These scenarios demand a communication protocol that
goes beyond simple request/response—one that handles partial failures
gracefully, reconciles mismatched capability interfaces, and preserves
human control over consequential actions.

Existing agent communication languages (FIPA-ACL, KQML) were designed
for structured knowledge sharing between software agents and predate
the era of LLM-backed agents that expose capabilities through
natural-language descriptions.  Contemporary frameworks such as
LangChain~\cite{harrison2022langchain}, Anthropic's Model Context
Protocol~\cite{anthropic2024mcp}, and Google's Agent-to-Agent
protocol~\cite{google2024a2a} address tool invocation but do not
provide transactional semantics for multi-step workflows or
first-class capability discovery.

The \emph{Aroha Communication Protocol}
(ACP)\footnote{Not to be confused with the \emph{Agentic Commerce
Protocol}, a 2025 checkout protocol for agent-initiated purchases.
Throughout this paper ACP refers exclusively to the Aroha
Communication Protocol; the two operate at different layers and are
unrelated.} fills this gap as the
semantic layer of the Aroha stack~\cite{aroha2025p5}.  Our
contributions are:
\begin{enumerate}
  \item A taxonomy of \textbf{seventeen message types} covering the full
        lifecycle of agent interactions (\S\ref{sec:messages}).
  \item \textbf{Four interaction patterns} with formal state machines,
        from two-message request/response to $4k$-message distributed
        sagas (\S\ref{sec:patterns}).
  \item A \textbf{Capability Schema Negotiation} (CSN) mechanism with
        three resolution paths and a persistent cache
        (\S\ref{sec:csn}).
  \item A \textbf{HITL approval gate} abstraction that plugs between
        saga reserve and commit phases (\S\ref{sec:hitl}).
  \item \textbf{Pluggable settlement} adapters for financial,
        quota-based, and on-chain escrow settlement
        (\S\ref{sec:settlement}).
  \item \textbf{Three formal theorems} with proofs: saga completeness,
        cancel idempotency, and CSN cache correctness
        (\S\ref{sec:formal}).
  \item \textbf{Satisfaction signals} that close the feedback loop
        to the Bayesian reputation system of Paper~1
        (\S\ref{sec:satisfaction}).
\end{enumerate}

ACP is implemented in the TypeScript packages
\texttt{@aroha-sdk/core} (message types, envelopes, validation) and
\texttt{@aroha-sdk/orchestrator} (saga engine, CSN), both published
on npm as part of the Aroha SDK.

% ============================================================
\section{Background and Related Work}
\label{sec:background}
% ============================================================

\subsection{Sagas}

Garcia-Molina and Salem~\cite{garcia1987sagas} introduced the saga as
a sequence of local transactions $T_1, T_2, \ldots, T_n$, each paired
with a compensating transaction $C_i$ that semantically undoes $T_i$.
A saga either completes all $T_i$ successfully, or executes $T_k$
partially and then compensates via $C_{k-1}, \ldots, C_1$ in reverse
order.  The key insight is that compensations are \emph{semantic}
undos executed as new transactions rather than rollbacks, making sagas
suitable for long-running workflows that span service boundaries.

ACP adapts this pattern to multi-agent networks by encoding
reservations as $T_i$ (via \texttt{ArohaReserve}/\texttt{ArohaReserveAck})
and cancellations as $C_i$ (via \texttt{ArohaCancel}/\texttt{ArohaCancelAck}).

\subsection{Capability Negotiation}

Schema-based service matching has been studied extensively in the
Web Services context (WSDL, OWL-S).  ACP's CSN mechanism extends this
with an LLM fallback for natural-language intent matching, inspired by
RAG-based document retrieval~\cite{lewis2020rag,karpukhin2020dpr}.
JSON Schema~\cite{w3cjsonschema} provides the structural type language;
SHA-256 fingerprints~\cite{jsoncanon} enable cache keying.

\subsection{Human-in-the-Loop Control}

HITL interfaces have been studied in the context of active
learning and supervised autonomy.  ACP's HITL gate is distinguished
by its \emph{transactional} positioning: it occupies the interval
between a reversible reservation and an irreversible commit, ensuring
that human approval is solicited only after a provider has confirmed
feasibility but before resources are consumed.

\subsection{Distributed Identity}

ACP messages carry Decentralized Identifier (DID)
references~\cite{w3c2022did} for agents and human credentials,
as specified in the companion protocol-stack paper~\cite{aroha2025p5}.
Ed25519 signatures~\cite{bernstein2011ed25519,josefsson2017ed25519rfc}
authenticate every message.

% ============================================================
\section{Message Type Taxonomy}
\label{sec:messages}
% ============================================================

ACP defines seventeen message types, organized into four groups.
Each message carries a standard envelope (version, messageId,
correlationId, sagaId, senderDID, recipientDID, timestamp, signature)
defined in~\cite{aroha2025p5}; Table~\ref{tab:messages} summarizes the
body fields specific to each type.

\subsection{Group 1: Direct Invocation}

\textbf{ArohaRequest} initiates a capability call.  Its body
carries: \texttt{capability} (string), \texttt{params} (JSON object),
optional \texttt{preferenceContext}, optional \texttt{idempotencyKey}
(24-hour deduplication window), and \texttt{credential}
(a \texttt{WithCredential} union carrying either a registry-based
\texttt{credentialId} or an inline base64url-encoded signed
\texttt{HumanCredential}).

\textbf{ArohaResponse} returns the capability result (\texttt{capability},
\texttt{result}).

\textbf{ArohaStream} delivers incremental progress: \texttt{capability},
\texttt{progressPct} $\in [0, 100]$, \texttt{event} string,
optional \texttt{data}.

\subsection{Group 2: Negotiation}

\textbf{ArohaNegotiate} proposes interaction terms:
\texttt{proposedTerms} $\{\textit{price}?,\ \textit{sla}?,\
\textit{validForMs}?\}$ plus an optional
\texttt{CSNNegotiateExtension} (see \S\ref{sec:csn}).

\textbf{ArohaCounterOffer} revises terms: \texttt{revisedTerms}
with the same schema, optional CSN extension.

\textbf{ArohaAccept} confirms agreement: \texttt{acceptedTerms},
optional \texttt{CSNAcceptExtension} carrying CSN resolution metadata.

\subsection{Group 3: Saga Execution}

\textbf{ArohaReserve} opens a saga step: \texttt{capability},
\texttt{params}, \texttt{holdDurationMs}, optional
\texttt{preferenceContext}, \texttt{credential}.

\textbf{ArohaReserveAck} confirms the reservation:
\texttt{reservationToken} (opaque string), \texttt{expiresAt} (ISO~8601).

\textbf{ArohaCommit} finalizes a step: \texttt{reservationToken}.

\textbf{ArohaCommitAck} delivers the step result:
\texttt{reservationToken}, \texttt{result}.

\textbf{ArohaCancel} compensates a step: \texttt{reservationToken},
\texttt{reason}.

\textbf{ArohaCancelAck} acknowledges cancellation:
\texttt{reservationToken}.

\subsection{Group 4: Meta}

\textbf{ArohaDelegate} forwards a request to another DID:
\texttt{targetDID}, \texttt{capability}, \texttt{params},
optional \texttt{preferenceContext}, \texttt{credential}.

\textbf{ArohaError} reports failures with a structured
\texttt{ArohaErrorCode} (17 values; Table~\ref{tab:errorcodes}),
\texttt{message}, \texttt{retryable} boolean,
optional \texttt{details} and \texttt{docsUrl}.

\textbf{ArohaSatisfactionSignal} closes the feedback loop
(see \S\ref{sec:satisfaction}).

\textbf{ArohaSpendingMandate} delegates a spending authority
(see \S\ref{sec:mandates}).

\textbf{ArohaTaskMandate} delegates non-monetary authority under
explicit resource constraints---allowed capabilities, per-capability
action limits, compute budgets, reversibility requirements, and
human-approval gates.  It mirrors the spending mandate's attenuation
semantics for task delegation.

\begin{table}[t]
\centering
\caption{ACP Message Type Taxonomy}
\label{tab:messages}
{\small
\begin{tabular}{@{}llp{3.0cm}@{}}
\toprule
\textbf{Group} & \textbf{Message Type} & \textbf{Key Body Fields} \\
\midrule
\multirow{3}{*}{Direct}
  & ArohaRequest      & capability, params, idempotencyKey \\
  & ArohaResponse     & capability, result \\
  & ArohaStream       & progressPct, event, data \\
\midrule
\multirow{3}{*}{Negotiation}
  & ArohaNegotiate    & proposedTerms, csn? \\
  & ArohaCounterOffer & revisedTerms, csn? \\
  & ArohaAccept       & acceptedTerms, csn? \\
\midrule
\multirow{6}{*}{Saga}
  & ArohaReserve      & capability, holdDurationMs \\
  & ArohaReserveAck   & reservationToken, expiresAt \\
  & ArohaCommit       & reservationToken \\
  & ArohaCommitAck    & reservationToken, result \\
  & ArohaCancel       & reservationToken, reason \\
  & ArohaCancelAck    & reservationToken \\
\midrule
\multirow{5}{*}{Meta}
  & ArohaDelegate          & targetDID, capability \\
  & ArohaError             & code, message, retryable \\
  & ArohaSatisfactionSignal & outcome, dimensions \\
  & ArohaSpendingMandate   & mandateTier, constraints \\
  & ArohaTaskMandate       & allowed, actionLimits, humanGates \\
\bottomrule
\end{tabular}
}
\end{table}

\begin{table}[t]
\centering
\caption{ArohaErrorCode Enumeration (17 values)}
\label{tab:errorcodes}
{\small
\begin{tabular}{@{}ll@{}}
\toprule
\textbf{Code} & \textbf{Meaning} \\
\midrule
Aroha\_UNAUTHORIZED         & Missing or invalid credential \\
Aroha\_FORBIDDEN            & Credential valid, access denied \\
Aroha\_INVALID\_SIGNATURE   & Envelope signature check failed \\
Aroha\_EXPIRED\_MESSAGE     & Message timestamp stale \\
Aroha\_REPLAY\_DETECTED     & Duplicate messageId seen \\
Aroha\_UNKNOWN\_CAPABILITY  & Capability not registered \\
Aroha\_TRUST\_LEVEL\_INSUFFICIENT & Insufficient reputation \\
Aroha\_RESERVATION\_FAILED  & Cannot lock resource \\
Aroha\_RESERVATION\_EXPIRED & Token TTL elapsed \\
Aroha\_COMMIT\_FAILED       & Commit-time settlement error \\
Aroha\_CANCEL\_FAILED       & Cancellation error (non-fatal) \\
Aroha\_INVALID\_TOKEN       & Malformed reservation token \\
Aroha\_NO\_INVENTORY        & Resource exhausted \\
Aroha\_CAPACITY\_EXCEEDED   & Rate/capacity limit hit \\
Aroha\_SERVICE\_UNAVAILABLE & Provider temporarily down \\
Aroha\_INTERNAL\_ERROR      & Unclassified provider error \\
Aroha\_INVALID\_PARAMS      & Parameter validation failure \\
CSN\_NO\_STRUCTURAL\_MATCH  & CSN match failed, LLM disallowed \\
\bottomrule
\end{tabular}
}
\end{table}

% ============================================================
\section{Interaction Patterns and Saga Lifecycle}
\label{sec:patterns}
% ============================================================

ACP specifies four canonical interaction patterns.  An orchestrator
selects a pattern based on whether transactional atomicity, streaming
progress, or capability discovery are required.

\subsection{Pattern 1: Simple (2 messages)}

\[
  \text{Orchestrator} \xrightarrow{\text{ArohaRequest}} \text{Provider}
  \xrightarrow{\text{ArohaResponse}|\text{ArohaError}} \text{Orchestrator}
\]

Suitable for idempotent read-like capabilities.  Total latency:
approximately one round-trip time (RTT), $\approx 5\,\text{ms}$ on
a local network.

\subsection{Pattern 2: Streamed ($2 + N$ messages)}

\[
  \text{ArohaRequest} \to
  \underbrace{\text{ArohaStream} \cdots \text{ArohaStream}}_{N}
  \to \text{ArohaResponse}
\]

Each \texttt{ArohaStream} carries a \texttt{progressPct} in $[0,100]$
and an event label.  The final \texttt{ArohaResponse} signals
completion.

\subsection{Pattern 3: Negotiated (4--6 messages)}

\[
  \text{Negotiate} \to [\text{CounterOffer} \to]^{?}
  \text{Accept} \to \text{Request} \to \text{Response}
\]

At most one counter-offer round is typical; the protocol imposes no
hard limit.  A \texttt{CSNNegotiateExtension} may be embedded to
perform capability discovery in the same round-trip as price
negotiation (see \S\ref{sec:csn}).

\subsection{Pattern 4: Distributed Saga ($4k$ messages)}

\begin{definition}[Aroha Saga]
An \emph{Aroha Saga} $\Saga = [(r_1, c_1), \ldots, (r_n, c_n)]$ is a
sequence of $n$ steps, where $r_i$ is the reserve-commit pair for step
$i$ and $c_i = \texttt{ArohaCancel}(\textit{token}_i)$ is the
corresponding compensating action.
\end{definition}

Algorithm~\ref{alg:saga} gives the orchestrator-side execution
procedure.  For each step $i$, the orchestrator sends
\texttt{ArohaReserve}, awaits \texttt{ArohaReserveAck} bearing an
opaque \texttt{reservationToken} $\tau_i$ and an expiry timestamp
$e_i$.  If a HITL gate is configured for this step, the orchestrator
suspends and polls for human approval before issuing
\texttt{ArohaCommit}($\tau_i$).  On receipt of
\texttt{ArohaCommitAck}($\tau_i, \textit{result}_i$), the step is
considered committed.

\textbf{Failure handling.}  If step $k$ fails (reserve fails, commit
fails, HITL rejects, or token expires), the saga enters the
\emph{compensation phase}: for $j = k, k-1, \ldots, 1$, the
orchestrator sends \texttt{ArohaCancel}($\tau_j$) and awaits
\texttt{ArohaCancelAck}($\tau_j$).  This LIFO order ensures
the most recently reserved resource is released first—e.g., a hotel
reservation is cancelled before the earlier flight booking it depended
on, matching the semantic dependency order.

\textbf{Idempotency.}  Every \texttt{ArohaRequest} and
\texttt{ArohaReserve} carries an optional \texttt{idempotencyKey}.
Providers deduplicate requests with the same key within a 24-hour
window, preventing duplicate side effects on network-level retries.

\begin{algorithm}[t]
\caption{Aroha Saga Execution (Orchestrator)}
\label{alg:saga}
{\small
\begin{algorithmic}[1]
\Require Saga $\Saga = [(r_1,c_1),\ldots,(r_n,c_n)]$, HITL gate $H$
\Ensure Saga terminates in SUCCESS or COMPENSATED
\State $\textit{tokens} \gets []$
\For{$i \gets 1$ \textbf{to} $n$}
  \State Send \texttt{ArohaReserve}($r_i$) to Provider$_i$
  \State $\textit{ack} \gets $ await \texttt{ArohaReserveAck} \textbf{or} \texttt{ArohaError}
  \If{$\textit{ack}$ is \texttt{ArohaError}}
    \State \textbf{goto} \textsc{Compensate}($i-1$)
  \EndIf
  \State $\tau_i \gets \textit{ack.reservationToken}$
  \State $\textit{tokens}.\text{append}(\tau_i)$
  \If{$H.\text{required}(\textit{ack})$}
    \State $d \gets H.\text{requestApproval}(\tau_i, r_i)$
    \If{$d = \text{``rejected''}$}
      \State \textbf{goto} \textsc{Compensate}($i$)
    \EndIf
  \EndIf
  \State Send \texttt{ArohaCommit}($\tau_i$) to Provider$_i$
  \State $\textit{res} \gets $ await \texttt{ArohaCommitAck} \textbf{or} \texttt{ArohaError}
  \If{$\textit{res}$ is \texttt{ArohaError}}
    \State \textbf{goto} \textsc{Compensate}($i$)
  \EndIf
\EndFor
\State \textbf{return} SUCCESS
\State
\Function{Compensate}{$k$}
  \For{$j \gets k$ \textbf{downto} $1$}
    \State Send \texttt{ArohaCancel}($\textit{tokens}[j]$) to Provider$_j$
    \State Await \texttt{ArohaCancelAck}($\textit{tokens}[j]$)
  \EndFor
  \State \textbf{return} COMPENSATED
\EndFunction
\end{algorithmic}
}
\end{algorithm}

\begin{table}[t]
\centering
\caption{Interaction Pattern Round-Trip Summary}
\label{tab:rtt}
{\small
\begin{tabular}{@{}llll@{}}
\toprule
\textbf{Pattern} & \textbf{Messages} & \textbf{RTTs} & \textbf{Est.\ Latency} \\
\midrule
Simple request      & 2       & 1       & $\approx 5\,\text{ms}$ \\
Streamed            & $2+N$   & $1+N$   & $5\,\text{ms}+N{\times}\Delta$ \\
Negotiated          & 4--6    & 2--3    & 10--15\,ms \\
Saga ($k$ steps)    & $4k$    & $2k$    & $10k\,\text{ms}$ \\
CSN (cache miss)    & 4--6    & 2--3    & 10--15\,ms + LLM \\
CSN (cache hit)     & 4       & 2       & $\approx 10\,\text{ms}$ \\
\bottomrule
\end{tabular}
}
\end{table}

% ============================================================
\section{Capability Schema Negotiation}
\label{sec:csn}
% ============================================================

Orchestrators in an open network may not know the precise capability
name or parameter schema expected by a provider.  CSN resolves this
mismatch through three paths, attempted in priority order.

\subsection{CSN Message Extensions}

A \texttt{CSNNegotiateExtension} is optionally embedded in
\texttt{ArohaNegotiate}; its fields are:
\begin{itemize}
  \item \texttt{capabilityIntent}: natural-language description
  \item \texttt{desiredSchema}: JSON Schema~\cite{w3cjsonschema} for
        input/output
  \item \texttt{schemaHash}: SHA-256 hex of canonical desiredSchema
  \item \texttt{requestCapabilityManifest}: boolean—ask provider to
        return its full capability manifest
  \item \texttt{disallowLLMFallback}: if true, force
        \texttt{CSN\_NO\_STRUCTURAL\_MATCH} rather than LLM inference
\end{itemize}

A \texttt{CSNAcceptExtension} is embedded in \texttt{ArohaAccept}:
\begin{itemize}
  \item \texttt{resolvedCapability}: canonical capability name
  \item \texttt{agreedSchema}: provider-normalized JSON Schema
  \item \texttt{schemaHash}: SHA-256 of agreedSchema
  \item \texttt{capabilityManifest}: full manifest if requested
  \item \texttt{matchMethod}: \texttt{"cache"} | \texttt{"schema"} |
        \texttt{"intent"}
  \item \texttt{compatibilityScore} $\in [0,1]$
  \item \texttt{matchAudit}: $\{\textit{structuralScore}?,
        \textit{cacheHit}?, \textit{llmUsed}, \textit{llmModel}?,
        \textit{fallbackReason}?\}$
\end{itemize}

\subsection{Resolution Path 1: Cache Hit}

If the orchestrator's \texttt{ArohaNegotiate} carries a
\texttt{schemaHash} $h$, the provider checks its
\texttt{IDiscoveryStore} at key $\textit{agentDID}\!:\!h$.  On a
cache hit, the provider returns the stored \texttt{capabilityName},
\texttt{agreedSchema}, and \texttt{matchMethod} in under 1\,ms
without LLM invocation.  The cache entry structure is:
\[
  \text{CacheEntry} = \{\textit{capabilityName},\ \textit{agreedSchema},\
  \textit{matchMethod},\ \textit{compatibilityScore}?,\
  \textit{cachedAt},\ \textit{hitCount}\}
\]

\subsection{Resolution Path 2: Schema-Based Matching}

If no cache entry exists but \texttt{desiredSchema} is provided,
the provider runs structural compatibility matching between the
requested schema and each of its registered capability schemas.
The compatibility score $\sigma \in [0,1]$ quantifies field overlap,
type compatibility, and required/optional field alignment.
If $\sigma \geq \theta$ (configurable threshold, default 0.85),
the match resolves without LLM invocation; \texttt{matchMethod}
is set to \texttt{"schema"}.

\subsection{Resolution Path 3: Intent-Based Matching}

If schema matching fails or \texttt{capabilityIntent} is the only
input, the provider invokes an LLM over its capability registry.
The LLM receives the intent string and a manifest of available
capabilities and returns the best match.  This path is skipped if
\texttt{disallowLLMFallback} is true.

\subsection{Post-Resolution Caching}

After any successful resolution, the provider stores the mapping:
\[
  \texttt{store.set}(\textit{agentDID}\!:\!h,\ \text{CacheEntry})
\]
The orchestrator similarly stores $(\textit{agentDID}, h,
\textit{capabilityName})$ locally.  All future negotiations to the
same provider with schema hash $h$ take Path~1.

\begin{algorithm}[t]
\caption{CSN Resolution (Provider-Side)}
\label{alg:csn}
{\small
\begin{algorithmic}[1]
\Require \texttt{CSNNegotiateExtension} $E$, agentDID $D$,
         capability registry $R$, store $\mathcal{C}$
\Ensure \texttt{CSNAcceptExtension} or \texttt{CSN\_NO\_STRUCTURAL\_MATCH}
\If{$E.\textit{schemaHash} \neq \texttt{null}$}
  \State $e \gets \mathcal{C}.\texttt{get}(D\!:\!E.\textit{schemaHash})$
  \If{$e \neq \texttt{null}$}
    \State $e.\textit{hitCount} \mathrel{+}= 1$
    \State \textbf{return} CSNAccept$(e,\ \texttt{matchMethod}=\texttt{"cache"})$
  \EndIf
\EndIf
\If{$E.\textit{desiredSchema} \neq \texttt{null}$}
  \State $(\hat{C}, \sigma) \gets \textsc{StructuralMatch}(E.\textit{desiredSchema}, R)$
  \If{$\sigma \geq \theta$}
    \State $h \gets \text{SHA-256}(\text{canonical}(E.\textit{desiredSchema}))$
    \State $\mathcal{C}.\texttt{set}(D\!:\!h,\ \{\hat{C}, E.\textit{desiredSchema}, \sigma\})$
    \State \textbf{return} CSNAccept$(\hat{C}, \sigma,\ \texttt{matchMethod}=\texttt{"schema"})$
  \EndIf
\EndIf
\If{$E.\textit{disallowLLMFallback} = \texttt{true}$}
  \State \textbf{throw} \texttt{CSN\_NO\_STRUCTURAL\_MATCH}
\EndIf
\If{$E.\textit{capabilityIntent} \neq \texttt{null}$}
  \State $\hat{C} \gets \textsc{LLMMatch}(E.\textit{capabilityIntent}, R)$
  \If{$\hat{C} \neq \texttt{null}$}
    \State \textbf{return} CSNAccept$(\hat{C},\ \texttt{matchMethod}=\texttt{"intent"})$
  \EndIf
\EndIf
\State \textbf{throw} \texttt{CSN\_NO\_STRUCTURAL\_MATCH}
\end{algorithmic}
}
\end{algorithm}

% ============================================================
\section{Human-in-the-Loop Approval Gate}
\label{sec:hitl}
% ============================================================

Many agent-driven workflows require human oversight before committing
irreversible or high-cost steps.  ACP provides a first-class HITL
abstraction that plugs between the reserve and commit phases of each
saga step.

\subsection{Approval Interface}

An approval gate implements a two-method interface:
\begin{align*}
  &\texttt{request}(\textit{sagaId}, \textit{stepId},
    \textit{HumanApprovalRequest}) \to \texttt{void} \\
  &\texttt{pollDecision}(\textit{sagaId}, \textit{stepId})
    \to \text{Decision} \mid \texttt{null}
\end{align*}
where \texttt{HumanApprovalRequest} carries \texttt{capability},
\texttt{agentDID}, \texttt{params}, optional \texttt{estimatedCost},
\texttt{costCurrency}, and \texttt{description}; and
$\text{Decision} \in \{\texttt{"approved"}, \texttt{"rejected"}\}$.

\subsection{Webhook Gate}

The built-in \texttt{WebhookApprovalGate} posts a
\texttt{HumanApprovalRequest} payload to a configured
\texttt{webhookUrl} over HTTPS, authenticating with an
\texttt{X-Aroha-Signature} HMAC header keyed by \texttt{secret}.
It then polls \texttt{pollDecision} at \texttt{pollIntervalMs}
(default 3\,000\,ms) until a decision arrives or
\texttt{timeoutMs} (default 300\,000\,ms = 5\,min) elapses.
On timeout, the gate returns \texttt{"rejected"}.

On \texttt{"rejected"}: the saga engine issues
\texttt{ArohaCancel}($\tau_i$) for the current step and initiates
LIFO compensation for all prior committed steps.

\subsection{Development Gate}

\texttt{AlwaysApproveGate} returns \texttt{"approved"} immediately and
is intended for testing environments only.

% ============================================================
\section{Settlement Adapters}
\label{sec:settlement}
% ============================================================

ACP decouples communication semantics from payment and resource
settlement through a pluggable \texttt{IArohaSettlement} interface.

\subsection{Settlement Interface}

{\small
\begin{align*}
  &\texttt{onReserve}(\textit{ctx}: \textit{ReserveCtx}) \to \texttt{void}
    \quad\text{[throw = reject step]} \\
  &\texttt{onCommit}(\textit{ctx}: \textit{CommitCtx}) \to \texttt{void}
    \quad\text{[throw = fail commit]} \\
  &\texttt{onCancel}(\textit{ctx}: \textit{CancelCtx}) \to \texttt{void}
    \quad\text{[errors logged, ignored]} \\
  &\texttt{onSagaComplete}(\textit{sagaId},\ \textit{outcome})
    \to \texttt{void}
\end{align*}
}
\noindent where $\textit{outcome} \in \{\texttt{"SUCCESS"},
\texttt{"COMPENSATED"}, \texttt{"FAILURE"}\}$.  Every context
carries: \texttt{sagaId}, \texttt{correlationId},
\texttt{orchestratorDID}, \texttt{providerDID}, \texttt{capability},
\texttt{params}.

\subsection{On-Chain Escrow}

The \texttt{OnChainSettlement} adapter delegates to
\texttt{IArohaEscrowClient}:
\begin{align*}
  &\texttt{open}(\textit{sagaId},\ \textit{ttlSeconds},\
    \textit{amountWei}) \to \textit{txHash} \\
  &\texttt{approveSplit}(\textit{sagaId},\ \textit{recipient},\
    \textit{amountWei}) \to \texttt{void} \\
  &\texttt{release}(\textit{sagaId}) \to \texttt{void}
    \quad\text{[saga succeeded]} \\
  &\texttt{refund}(\textit{sagaId}) \to \texttt{void}
    \quad\text{[saga compensated]}
\end{align*}
\texttt{onReserve} calls \texttt{open}; \texttt{onCommit} calls
\texttt{approveSplit}; \texttt{onSagaComplete(SUCCESS)} calls
\texttt{release}; \texttt{onSagaComplete(COMPENSATED)} calls
\texttt{refund}.  The contract~\cite{buterin2014ethereum} enforces
atomicity on-chain.

\subsection{Available Implementations}

Table~\ref{tab:settlement} summarizes the five built-in adapters.

\begin{table}[t]
\centering
\caption{Built-in Settlement Adapter Implementations}
\label{tab:settlement}
{\small
\begin{tabular}{@{}lp{4.5cm}@{}}
\toprule
\textbf{Adapter} & \textbf{Description} \\
\midrule
\texttt{StripeSettlement}    & Locks via payment intents; captures
                               on commit; refunds on compensation \\
\texttt{PostgresSettlement}  & Atomic ledger deductions in a SQL
                               transaction; release on cancel \\
\texttt{OnChainSettlement}   & EVM escrow contract; multi-party splits \\
\texttt{QuotaSettlement}     & Decrements API-key quota counter;
                               restores on cancel \\
\texttt{NullSettlement}      & No-op; for capabilities with no
                               financial settlement \\
\bottomrule
\end{tabular}
}
\end{table}

% ============================================================
\section{Spending Mandates}
\label{sec:mandates}
% ============================================================

When an orchestrator delegates to a sub-agent, it issues an
\texttt{ArohaSpendingMandate} granting bounded spending authority.
This complements the mandate-chain mechanism of Paper~2~\cite{aroha2025p2}.

\subsection{Mandate Structure}

Each mandate carries:
\begin{itemize}
  \item \texttt{mandateTier} $\in \{\texttt{"intent"},
        \texttt{"cart"}, \texttt{"payment"}\}$
  \item \texttt{constraints}: a \texttt{SpendingConstraints} object
        with \texttt{spendLimitUsd}, optional \texttt{sessionLimitUsd},
        \texttt{currency}, \texttt{merchantCategory},
        \texttt{requireHumanApprovalAboveUsd},
        \texttt{allowedMerchants}, \texttt{validFrom}, \texttt{validUntil}
  \item \texttt{grantee}, \texttt{grantor}: DIDs
  \item \texttt{issuedAt}, \texttt{expiresAt}: ISO~8601 timestamps
  \item \texttt{mandateId}, \texttt{parentMandateId}: chain linkage
\end{itemize}

\textbf{Mandate tiers.}  An \emph{intent} mandate authorizes planning
and price discovery only.  A \emph{cart} mandate authorizes reservation
(reserve phase) but not commit.  A \emph{payment} mandate authorizes
full saga execution including commit and financial settlement.

\subsection{Dual-Layer Budget Enforcement}

Correct delegation requires two enforcement layers acting
in concert.

\textbf{Layer 1: Monotonic narrowing.}  At attenuation time, the
system checks:
\[
  \textit{child.spendLimitUsd} \leq \textit{parent.spendLimitUsd}
\]
This is verified synchronously before the child mandate is signed.
Narrowing alone is insufficient: $N$ children can each claim the
full parent limit, yielding an effective grant of
$N \times \textit{spendLimitUsd}$.

\textbf{Layer 2: Cumulative ledger.}  A shared atomic counter tracks
the total budget allocated across all child mandates of a given parent.
When a child mandate is signed, the ledger atomically deducts
\texttt{spendLimitUsd} from the parent pool.  If the deduction would
exceed the remaining balance, a \texttt{BudgetExhaustedError} is
thrown before signing.  When a saga is compensated, the unused
allocation is returned via \texttt{releaseAllocation()}.

\textbf{Why both layers are necessary.}  Narrowing without a ledger
permits over-allocation.  A ledger without narrowing permits a child
to re-widen a narrowed constraint.  Together they establish:
\[
  \sum_j \textit{child}_j.\textit{spendLimitUsd}
  \leq \textit{parent.spendLimitUsd}
  \quad \text{and} \quad
  \textit{child}_j \subseteq \textit{parent} \;\forall j
\]
ensuring total allocation is bounded and every child is a strict
subset of its parent's authority.

\begin{table}[t]
\centering
\caption{Spending Mandate Tier Semantics}
\label{tab:mandates}
{\small
\begin{tabular}{@{}lp{4.5cm}@{}}
\toprule
\textbf{Tier} & \textbf{Authorized Actions} \\
\midrule
\texttt{intent}  & Price discovery, negotiation only; no reservations \\
\texttt{cart}    & ArohaReserve permitted; ArohaCommit forbidden \\
\texttt{payment} & Full saga execution; financial settlement authorized \\
\bottomrule
\end{tabular}
}
\end{table}

% ============================================================
\section{Human Credentials and Roles}
\label{sec:credentials}
% ============================================================

ACP messages that carry a human credential use the
\texttt{WithCredential} interface:
\begin{itemize}
  \item \texttt{credentialId}: a registry-resolvable short identifier
        (preferred; enables revocation)
  \item \texttt{credentialToken}: inline base64url-encoded signed
        \texttt{HumanCredential} (used on first contact or absent
        shared registry)
\end{itemize}
A \texttt{HumanCredential} carries: \texttt{credentialId},
\texttt{userId}, optional \texttt{email}, \texttt{roles[]} (an array
of \texttt{ArohaRole} values), \texttt{issuedAt}, \texttt{expiresAt},
\texttt{issuerDID}, and an Ed25519 \texttt{signature}~\cite{bernstein2011ed25519}.

\textbf{ArohaRole.}  Six roles form a two-namespace hierarchy:
\texttt{human:admin}, \texttt{human:user}, \texttt{human:readonly}
for human principals; and \texttt{agent:orchestrator},
\texttt{agent:provider}, \texttt{agent:observer} for agent principals.
Role presence is checked at the provider before saga steps are accepted.

% ============================================================
\section{Formal Properties}
\label{sec:formal}
% ============================================================

We prove three key properties of ACP.

\begin{theorem}[Saga Completeness]
\label{thm:saga}
For any Aroha Saga $\Saga = [(r_1, c_1), \ldots, (r_n, c_n)]$,
every terminal state of saga execution is either:
\begin{enumerate}
  \item \textbf{SUCCESS}: all steps $r_1, \ldots, r_n$ are committed; or
  \item \textbf{COMPENSATED}: for some $k \in \{1,\ldots,n\}$,
        steps $r_1, \ldots, r_{k-1}$ are committed and compensations
        $c_{k-1}, \ldots, c_1$ have been executed in LIFO order.
\end{enumerate}
No step is left in a \emph{reserved-but-not-resolved} state.
\end{theorem}

\begin{proof}
By induction on $n$.

\textbf{Base case} ($n = 1$): \texttt{ArohaReserve} either returns
\texttt{ArohaReserveAck} or \texttt{ArohaError}.  In the error case,
no token was issued, so no compensation is needed; the saga
terminates in COMPENSATED (vacuously, no $c_j$ to execute).
If \texttt{ArohaReserveAck} is returned, either the optional HITL
gate approves and \texttt{ArohaCommit} is sent, or the gate rejects
and \texttt{ArohaCancel}($\tau_1$) is sent.  In the commit branch,
\texttt{ArohaCommitAck} (SUCCESS) or \texttt{ArohaError}
(triggers \texttt{ArohaCancel}($\tau_1$)) follows.  In all cases,
$\tau_1$ is either committed or cancelled; no reserved-but-unresolved
state is reachable.

\textbf{Inductive step}: Assume the property holds for all sagas of
length $< n$.  Consider $\Saga$ of length $n$.  Steps $1, \ldots, k-1$
succeed by assumption.  At step $k$, any failure (reserve error, commit
error, HITL rejection, token expiry) triggers Algorithm~\ref{alg:saga}'s
\textsc{Compensate}($k$) subroutine, which issues
\texttt{ArohaCancel}($\tau_j$) for $j = k, k-1, \ldots, 1$.
The \textsc{Compensate} loop is bounded ($k \leq n$) and terminates.
By Theorem~\ref{thm:cancel} (proved below), each \texttt{ArohaCancel}
is idempotent, so retries on transient network failures do not alter
the final state.  Thus every step is either committed or compensated;
the reserved-but-unresolved state is unreachable for $\Saga$.
\end{proof}

\begin{theorem}[ArohaCancel Idempotency]
\label{thm:cancel}
For any reservation token $\tau$, executing $\texttt{ArohaCancel}(\tau)$
one or more times produces the same final state as executing it exactly
once.
\end{theorem}

\begin{proof}
On the \emph{first} receipt of $\texttt{ArohaCancel}(\tau)$, the
provider atomically transitions the reservation with token $\tau$ from
state \textsc{Reserved} to state \textsc{Cancelled}, records the
cancellation, and returns $\texttt{ArohaCancelAck}(\tau)$.

On \emph{subsequent} receipt of $\texttt{ArohaCancel}(\tau)$, the
provider finds the reservation already in state \textsc{Cancelled}.
It returns $\texttt{ArohaCancelAck}(\tau)$ without modifying any
further state.

Since the state transitions are: \textsc{Reserved} $\to$
\textsc{Cancelled} (first call), and \textsc{Cancelled} $\to$
\textsc{Cancelled} (idempotent thereafter), the final state is
identical regardless of repetition count.
\end{proof}

\begin{theorem}[CSN Cache Correctness]
\label{thm:csn}
Let schema $S$ have SHA-256 hash $h = \text{SHA-256}(\text{canonical}(S))$.
If a CSN resolution for $(D, h)$ (provider DID $D$, hash $h$) has been
stored in the \texttt{IDiscoveryStore} with capability name $\hat{C}$
and agreed schema $A$, then for all future \texttt{ArohaNegotiate}
messages from any orchestrator carrying \texttt{schemaHash} $= h$
directed to $D$, the \texttt{ArohaAccept} response carries
\texttt{resolvedCapability} $= \hat{C}$ and \texttt{agreedSchema} $= A$
without LLM invocation.
\end{theorem}

\begin{proof}
The cache is keyed by $D\!:\!h$ and set atomically on successful
resolution (Algorithm~\ref{alg:csn}, lines 12--13 and 19--20).
The key $D\!:\!h$ is deterministic given $D$ and $h$.

For correctness, we require that two distinct schemas $S_1 \neq S_2$
cannot share hash $h$.  This follows from the collision resistance of
SHA-256: an adversary producing $S_1 \neq S_2$ with
$\text{SHA-256}(\text{canonical}(S_1)) = \text{SHA-256}(\text{canonical}(S_2))$
would constitute a collision attack with probability bounded by
$2^{-128}$ under standard cryptographic assumptions~\cite{keccak2015}.

Given collision resistance, the cache is injective: at most one entry
exists per key $D\!:\!h$, and its content is fixed at write time.
Every subsequent \texttt{store.get}($D\!:\!h$) returns the same
$(\hat{C}, A)$.  Algorithm~\ref{alg:csn} returns the cache result
immediately at lines 4--7, bypassing all schema-matching and LLM
paths.
\end{proof}

% ============================================================
\section{Satisfaction Signals and Reputation Feedback}
\label{sec:satisfaction}
% ============================================================

ACP closes the quality feedback loop between saga outcomes and the
Bayesian reputation system of Paper~1~\cite{aroha2025p1} through
the \texttt{ArohaSatisfactionSignal} message type.

\subsection{Signal Structure}

\texttt{ArohaSatisfactionSignalBody} carries:
\begin{itemize}
  \item \texttt{sagaId}, \texttt{agentDID}, \texttt{capability}
  \item \texttt{outcome} $\in \{\texttt{"satisfied"},
        \texttt{"neutral"}, \texttt{"dissatisfied"}\}$
  \item \texttt{dimensions}: $\{\textit{price}, \textit{quality},
        \textit{speed}\}$, each a 3-way signal
  \item \texttt{withinSLA}: boolean
  \item \texttt{requiredCompensation}: boolean (true if a saga
        compensation was triggered)
\end{itemize}

\subsection{Reputation Mapping}

Satisfaction signals map to the $\beta$-distribution parameters
of Paper~1's \texttt{ArohaRep} system~\cite{josang2002beta} as follows:
\begin{align*}
  \texttt{outcome}&=\texttt{"satisfied"} \wedge
    \texttt{withinSLA}=\texttt{true}
    &\Rightarrow&\quad s^{+}\ (\text{weight}\ 1.0) \\
  \texttt{outcome}&=\texttt{"dissatisfied"}
    &\Rightarrow&\quad s^{-}\ (\text{weight}\ 1.2) \\
  \texttt{requiredCompensation}&=\texttt{true}
    &\Rightarrow&\quad s^{-}\ (\text{weight}\ 0.8)
\end{align*}

The \texttt{IReputationRecorder} interface accepts a
\texttt{ReputationRecordContext} carrying: \texttt{sagaId},
\texttt{agentDID}, \texttt{outcome}, optional
\texttt{latencyMs}, \texttt{slaDurationMs}, \texttt{quotedPrice},
\texttt{actualPrice}, \texttt{retryCount}, \texttt{qualityScore},
\texttt{domain}.  The reputation system uses these fields to maintain
separate per-domain and per-capability reputation models.

\textbf{Compensation as quality signal.}  A saga that requires
LIFO compensation indicates partial failure—even if the orchestrator's
perspective is neutral (no harm done), the provider failed to complete
a committed step.  Flagging \texttt{requiredCompensation = true}
allows the reputation system to penalize providers whose workflows
require frequent rollbacks.

% ============================================================
\section{Distributed Tracing and Observability}
\label{sec:tracing}
% ============================================================

ACP messages carry a W3C Trace Context~\cite{w3c2021tracecontext}
\texttt{correlationId} in the envelope, enabling cross-agent
distributed tracing.  The \texttt{sagaId} provides coarser-grained
grouping of all messages belonging to a single saga execution.

Observability spans are emitted at each saga phase transition:
Reserve, CommitPending, Committed, Cancelled.  Implementations
are expected to export spans to OpenTelemetry-compatible collectors.
The \texttt{matchAudit} field of \texttt{CSNAcceptExtension} provides
transparency into which CSN resolution path was taken and whether an
LLM was invoked, supporting cost accounting and compliance auditing.

% ============================================================
\section{Security Considerations}
\label{sec:security}
% ============================================================

\textbf{Replay attacks.}  ACP envelope timestamps are checked against
a configurable maximum clock skew.  The \texttt{messageId} field
is stored in a sliding-window cache; duplicates return
\texttt{Aroha\_REPLAY\_DETECTED}.  The \texttt{idempotencyKey}
provides higher-level deduplication at the business logic layer.

\textbf{Credential theft.}  Inline \texttt{credentialToken} values
are base64url-encoded signed JWTs; their \texttt{expiresAt} field
limits exposure.  Registry-based \texttt{credentialId} references
can be revoked immediately by the issuer.

\textbf{HITL injection.}  Webhook payloads are authenticated with
HMAC-SHA256 (\texttt{X-Aroha-Signature}); receivers must verify
the signature before processing approval decisions.

\textbf{CSN adversarial schemas.}  An adversary could supply a
crafted \texttt{desiredSchema} to trigger expensive LLM inference
on every request.  The \texttt{schemaHash} cache and
\texttt{disallowLLMFallback} flag mitigate this.  Rate limiting
on the \texttt{Intent} path is recommended for public endpoints.

\textbf{Mandate over-allocation.}  Prevented by the dual-layer
enforcement described in \S\ref{sec:mandates}: monotonic narrowing
plus atomic ledger deduction.

% ============================================================
\section{Implementation and Evaluation}
\label{sec:impl}
% ============================================================

ACP is implemented as the TypeScript package
\texttt{@aroha-labs/acp}, published on npm.  The package exports:
the seventeen message constructors and validators, the saga engine
(Algorithm~\ref{alg:saga}), the CSN resolver
(Algorithm~\ref{alg:csn}), \texttt{WebhookApprovalGate} and
\texttt{AlwaysApproveGate}, the five settlement adapters, and the
\texttt{IDiscoveryStore} interface with Redis, Postgres, and
in-memory backends.

\subsection{Micro-Benchmarks}

Table~\ref{tab:bench} reports mean latencies measured on a
single host (AMD Ryzen 9, 64\,GB RAM, Node.js 22) for 10\,000
iterations.

\begin{table}[t]
\centering
\caption{ACP Operation Micro-Benchmarks (single-host)}
\label{tab:bench}
{\small
\begin{tabular}{@{}lrr@{}}
\toprule
\textbf{Operation} & \textbf{Mean} & \textbf{P99} \\
\midrule
Envelope sign (Ed25519)        & 0.08\,ms & 0.12\,ms \\
Envelope verify (Ed25519)      & 0.06\,ms & 0.09\,ms \\
Schema hash (SHA-256, 4\,KB)   & 0.02\,ms & 0.04\,ms \\
CSN cache lookup (Redis)       & 0.9\,ms  & 1.4\,ms  \\
CSN structural match (10 caps) & 2.1\,ms  & 3.8\,ms  \\
CSN intent (LLM, GPT-4o-mini)  & 480\,ms  & 920\,ms  \\
Saga step (reserve+commit)     & 4.8\,ms  & 9.2\,ms  \\
\bottomrule
\end{tabular}
}
\end{table}

The data confirm that CSN cache hits ($< 1\,\text{ms}$) and
structural matching ($< 4\,\text{ms}$) impose negligible overhead
on a per-invocation basis.  LLM-assisted intent matching is
appropriate only for initial discovery; the cache ensures subsequent
interactions avoid it.

% ============================================================
\section{Discussion}
\label{sec:discussion}
% ============================================================

\textbf{Comparison with database sagas.}  Classic sagas assume
shared database infrastructure for coordination~\cite{garcia1987sagas}.
ACP's saga pattern replaces this with message-passing over the
Aroha envelope protocol~\cite{aroha2025p5}, which carries
distributed identity and cryptographic authentication.
The compensation model is analogous: ArohaCancel is a semantic undo,
not a database rollback.  The LIFO ordering preserves the
``undo the most recent action first'' invariant that Garcia-Molina
and Salem identified as essential for semantic consistency.

\textbf{Comparison with choreography vs.\ orchestration.}  ACP's
saga pattern is orchestration-centric: a single orchestrator
coordinates all reserve/commit/cancel messages.  Choreography-based
approaches distribute this responsibility, reducing single points of
failure but complicating LIFO ordering.  ACP's design prioritizes
correctness (Theorem~\ref{thm:saga}) and auditability (single
trace ID per saga) over decentralization at the coordination level.
Agent identity remains fully decentralized via DIDs~\cite{w3c2022did}.

\textbf{CSN and capability discovery.}  CSN's three-path design
mirrors the precision/recall trade-off in information retrieval:
schema matching provides high precision at low cost; intent-based LLM
matching provides higher recall at higher cost.  The cache converts
the one-time LLM cost into amortized near-zero cost for repeat
interactions~\cite{lewis2020rag}.  This is particularly valuable in
long-running agent deployments where the same orchestrator repeatedly
invokes the same provider.

\textbf{HITL gate placement.}  Positioning the HITL gate between
reserve and commit is a deliberate design choice.  A pre-reserve gate
requires the human to approve based on incomplete information (the
provider has not yet confirmed feasibility or price).  A post-commit
gate provides no opportunity to prevent the action.  The
reserve/commit split provides the ``reversible hold'' window that
makes meaningful human oversight practical.

% ============================================================
\section{Conclusion}
\label{sec:conclusion}
% ============================================================

We presented ACP, the semantic communication layer of the Aroha
agent network stack.  ACP's seventeen message types, four interaction
patterns, saga execution engine, and Capability Schema Negotiation
mechanism provide a complete foundation for multi-agent workflows
that require transactional atomicity, capability discovery, human
oversight, and financial settlement.

Three formal theorems establish that: sagas always terminate cleanly
without orphaned reservations; cancellations are safe to retry; and
CSN cache entries are correct by the collision resistance of SHA-256.

ACP is designed to complement the Aroha protocol stack~\cite{aroha2025p5},
reputation system~\cite{aroha2025p1}, mandate-chain
authorization~\cite{aroha2025p2}, publisher verification~\cite{aroha2025p3},
and adaptive RAG~\cite{aroha2025p4} of the companion papers.
Together they form a principled, open specification for deploying
trustworthy autonomous AI agent networks at scale.

\bibliographystyle{IEEEtran}
\bibliography{aroha}

\end{document}
